% ===== PRÉAMBULE CANONIQUE — à coller VERBATIM en tête de CHAQUE .tex =====
\documentclass[11pt,a4paper]{article}
\usepackage[utf8]{inputenc}\usepackage[T1]{fontenc}\usepackage{lmodern}
\usepackage[french]{babel}
\usepackage{amsmath,amssymb,mathtools}\usepackage{icomma}
\usepackage[version=4]{mhchem}          % \ce{...} : équations & formules
\usepackage{siunitx}\sisetup{locale=FR} % \qty{}{}, \si{}, \ang{}
\usepackage{chemfig}                    % \chemfig{...} : structures développées
\usepackage{xcolor}\usepackage{colortbl}
\usepackage{tikz}\usepackage{pgfplots}\pgfplotsset{compat=1.18}
\usepackage[most]{tcolorbox}\usepackage{enumitem}\usepackage{array,booktabs}
\usepackage[a4paper,margin=2.2cm]{geometry}
\usetikzlibrary{arrows.meta,calc,angles,quotes,positioning,patterns,backgrounds,shapes.geometric}
\definecolor{primary}{RGB}{222,49,99}\definecolor{primarydark}{RGB}{150,22,55}
\definecolor{defcol}{RGB}{0,114,178}\definecolor{methcol}{RGB}{52,92,148}
\definecolor{excol}{RGB}{0,135,90}\definecolor{attncol}{RGB}{200,40,55}
\tcbset{boxbase/.style={enhanced,breakable,boxrule=0.9pt,arc=2.5pt,
  left=3mm,right=3mm,top=2.3mm,bottom=2.3mm,fonttitle=\bfseries,coltitle=white}}
% --- boîtes TUTORIEL (noms EXACTS, ne pas renommer) ---
\newtcolorbox{cle}[1][]{boxbase,colback=defcol!6,colframe=defcol,
  title={\textsc{À retenir}\ifx\relax#1\relax\relax\else\textmd{ : \itshape #1}\fi}}
\newtcolorbox{methode}[1][]{boxbase,colback=methcol!7,colframe=methcol,sharp corners,
  title={\textsc{Méthode}\ifx\relax#1\relax\relax\else\textmd{ --- \itshape #1}\fi}}
\newtcolorbox{exemplebox}[1][]{boxbase,colback=excol!5,colframe=excol,
  title={\textsc{Exemple}\ifx\relax#1\relax\relax\else\textmd{ : \itshape #1}\fi}}
\newtcolorbox{piege}[1][]{boxbase,colback=attncol!6,colframe=attncol,
  title={\textsc{Piège}\ifx\relax#1\relax\relax\else\textmd{ : \itshape #1}\fi}}
\newtcolorbox{remarque}[1][]{boxbase,colback=black!4,colframe=black!45,
  title={\textsc{Remarque}\ifx\relax#1\relax\relax\else\textmd{ : \itshape #1}\fi}}
% --- boîtes EXERCICE / CORRIGÉ (noms EXACTS, auto-comptées) ---
\newtcolorbox[auto counter]{exo}[1][]{boxbase,colback=primary!4,colframe=primary,
  title={\textsc{Exercice}~\thetcbcounter\ifx\relax#1\relax\relax\else\textmd{ --- \itshape #1}\fi}}
\newtcolorbox[auto counter]{corrige}[1][]{boxbase,colback=excol!4,colframe=excol,
  title={\textsc{Corrigé}~\thetcbcounter\ifx\relax#1\relax\relax\else\textmd{ --- \itshape #1}\fi}}
\newcommand{\R}{\mathbb{R}}\newcommand{\N}{\mathbb{N}}\newcommand{\Z}{\mathbb{Z}}\newcommand{\ds}{\displaystyle}
% (un \usepackage{fancyhdr}+en-tête est possible ; il est ignoré côté web)
% =========================================================================
\begin{document}
\begin{corrige}[capacité des couches]
\begin{enumerate}[label=\alph*)]
  \item K : $2\times 1^2 = 2$ ; L : $2\times 2^2 = 8$ ; M : $2\times 3^2 = 18$ électrons.
  \item K $\leftrightarrow n=1$, L $\leftrightarrow n=2$, M $\leftrightarrow n=3$.
  \item N : $2\times 4^2 = 32$ électrons.
\end{enumerate}
\end{corrige}

\begin{corrige}[configuration d'atomes légers]
On remplit K ($2$), puis L ($8$), puis M.
\begin{enumerate}[label=\alph*)]
  \item $\ce{C}$ : $6$ électrons $\to (\mathrm{K})^2(\mathrm{L})^4$.
  \item $\ce{Na}$ : $11$ électrons $\to (\mathrm{K})^2(\mathrm{L})^8(\mathrm{M})^1$.
  \item $\ce{Cl}$ : $17$ électrons $\to (\mathrm{K})^2(\mathrm{L})^8(\mathrm{M})^7$.
\end{enumerate}
\end{corrige}

\begin{corrige}[couche de valence et électrons de valence]
La couche de valence est la dernière occupée.
\begin{enumerate}[label=\alph*)]
  \item $\ce{O}$ : $(\mathrm{K})^2(\mathrm{L})^6$ ; valence L, \textbf{$6$} électrons de valence.
  \item $\ce{Mg}$ : $(\mathrm{K})^2(\mathrm{L})^8(\mathrm{M})^2$ ; valence M, \textbf{$2$} électrons de valence.
  \item $\ce{Ar}$ : $(\mathrm{K})^2(\mathrm{L})^8(\mathrm{M})^8$ ; valence M, \textbf{$8$} électrons de
        valence (octet : gaz noble).
\end{enumerate}
\end{corrige}

\begin{corrige}[remonter à l'élément]
$Z$ = somme des exposants.
\begin{enumerate}[label=\alph*)]
  \item $2+5 = 7 \Rightarrow Z = 7$ : azote $\ce{N}$.
  \item $2+8+3 = 13 \Rightarrow Z = 13$ : aluminium $\ce{Al}$.
  \item $2+8+8 = 18 \Rightarrow Z = 18$ : argon $\ce{Ar}$.
\end{enumerate}
\end{corrige}

\begin{corrige}[symbole de Lewis]
On place les électrons de valence un par côté, puis on apparie ; les paires sont des doublets, les
points seuls des célibataires.
\begin{enumerate}[label=\alph*)]
  \item $\ce{C}$ ($4$) : $0$ doublet, $\mathbf{4}$ célibataires.
  \item $\ce{N}$ ($5$) : $\mathbf{1}$ doublet, $\mathbf{3}$ célibataires.
  \item $\ce{F}$ ($7$) : $\mathbf{3}$ doublets, $\mathbf{1}$ célibataire.
\end{enumerate}

\begin{center}
\begin{tikzpicture}[baseline]
  % --- Carbone : 4 célibataires ---
  \node[font=\Large] (C) at (0,0) {\textbf{C}};
  \fill (0,0.5) circle (1.5pt); \fill (0,-0.5) circle (1.5pt);
  \fill (-0.5,0) circle (1.5pt); \fill (0.5,0) circle (1.5pt);
  \node[font=\footnotesize] at (0,-1.15) {C};
  % --- Azote : 1 doublet + 3 célibataires ---
  \node[font=\Large] (N) at (3,0) {\textbf{N}};
  \fill (2.88,0.5) circle (1.5pt); \fill (3.12,0.5) circle (1.5pt); % doublet haut
  \fill (3,-0.5) circle (1.5pt);
  \fill (2.5,0) circle (1.5pt); \fill (3.5,0) circle (1.5pt);
  \node[font=\footnotesize] at (3,-1.15) {N};
  % --- Fluor : 3 doublets + 1 célibataire ---
  \node[font=\Large] (F) at (6,0) {\textbf{F}};
  \fill (5.88,0.5) circle (1.5pt); \fill (6.12,0.5) circle (1.5pt);   % doublet haut
  \fill (5.88,-0.5) circle (1.5pt); \fill (6.12,-0.5) circle (1.5pt); % doublet bas
  \fill (5.5,-0.12) circle (1.5pt); \fill (5.5,0.12) circle (1.5pt);  % doublet gauche
  \fill (6.5,0) circle (1.5pt);                                       % célibataire droite
  \node[font=\footnotesize] at (6,-1.15) {F};
\end{tikzpicture}
\end{center}
\end{corrige}

\end{document}
